% titel: Blatt 6, A7: VSEPR – Struktur, Winkel, Hybridisierung

\subsection*{Lösung}
\begin{center}
\begin{tabular}{@{}l c c c@{}}
  Molekül & EP-Anordnung (Zentralatom) & Geometrie \& Winkel & Hybridisierung \\
  \hline
  $\mathrm{H_2O}$ & tetraedrisch (4 EP, 2 frei) & gewinkelt (bent),
      $104.5^\circ$ & $sp^3$ \\
  $\mathrm{NH_3}$ & tetraedrisch (4 EP, 1 frei) & trigonal-pyramidal,
      $\approx 107^\circ$ & $sp^3$ \\
  $\mathrm{CO_2}$ & linear (2 EP)               & linear, $180^\circ$       & $sp$  \\
  $\mathrm{BH_3}$ & trigonal-planar (3 EP)      & trigonal-planar,
      $120^\circ$    & $sp^2$ \\
\end{tabular}
\end{center}

\subsection*{Begründung}
\begin{itemize}
  \item \textbf{$\mathrm{H_2O}$:} O bringt 2 $\sigma$-Bindungen und 2 freie
        Paare $\to$ 4 EP-Gebiete, $sp^3$; die zwei freien Paare drücken die
        Bindung auf $104.5^\circ$ (ideal $109.5^\circ$).
  \item \textbf{$\mathrm{NH_3}$:} N 3 Bindungen + 1 freies Paar $\to$ 4
        EP-Gebiete, $sp^3$; das eine freie Paar komprimiert auf $\approx 107^\circ$.
  \item \textbf{$\mathrm{CO_2}$:} C 2 Doppelbindungen $\to$ 2 EP-Gebiete,
        linear, $sp$.
  \item \textbf{$\mathrm{BH_3}$:} B 3 Bindungen, kein freies Paar $\to$ 3
        EP-Gebiete, trigonal-planar, $sp^2$, $120^\circ$.
\end{itemize}

\begin{ergebnis}
$\mathrm{H_2O}$: gewinkelt/$104.5^\circ$/$sp^3$; $\mathrm{NH_3}$:
trigonal-pyramidal/$\approx 107^\circ$/$sp^3$; $\mathrm{CO_2}$:
linear/$180^\circ$/$sp$; $\mathrm{BH_3}$: trigonal-planar/$120^\circ$/$sp^2$.
\end{ergebnis}

\begin{kontrolle}
Die EP-Zahlen (4, 4, 2, 3) und die dazugehörigen Hybridisierungen
($sp^3, sp^3, sp, sp^2$) sind konsistent; der Winkel-``Zug'' durch freie Paare
stimmt (Wasser am stärksten, dann NH$_3$). Beobachtete Werte ($104.5^\circ$,
$107^\circ$, $180^\circ$, $120^\circ$) bestätigen die Vorhersage.
\end{kontrolle}

\begin{merksatz}
EP-Paare zählen (Bindungen + freie Paare), Anordnung und Hybridisierung
ablesen; freie Paare drücken die Winkel zusammen. Doppelbindungen zählen als ein
EP-Gebiet.
\end{merksatz}
